% PDF 150、152--160；原书第 120、122--130 页。
\section{教程一——比较文本分类与命名实体识别中的提示和预训练再微调方法}
\label{sec:ch03-tutorial-one}
\subsection{概述}
本章介绍了基于提示的学习的概念，并详细讨论了提示与答案形式的几种配置。与 PTFT 相比，基于提示的方法的主要优点之一，是它在零样本和少样本条件下的表现。这一点还没有演示过。本教程将说明，与传统的基于任务头的微调相比，基于提示的学习怎样用更少的训练样例取得更好的结果。借助这一特点，可以用较少的数据和算力让大语言模型适应新任务。

\begin{keypoint}
\noindent\textbf{目标}
\begin{itemize}
\item 比较基于提示的学习与基于任务头的微调。
\item 演示怎样组织提示以完成不同的任务。
\item 介绍 OpenPrompt 库，说明本章介绍的技术是怎样实现的。
\item 绘制学习曲线，展示提示在少样本设置下的表现。
\end{itemize}
\end{keypoint}

这里给出教程的简明版本。完整版见 \url{https://github.com/springer-llms-deep-dive/llms-deep-dive-tutorials}。

\subsection{实验设计}
实验将直接比较 PTFT 和基于提示的学习在零样本及少样本条件下做文本分类和命名实体识别的能力。两种微调实验都以 BERT 为基础，使用 PyTorch，基于提示的部分还使用 OpenPrompt。逐步增加训练数据子集的大小，依次调优 BERT 模型，并定期在验证集上做预测，观察模型在少样本条件下怎样学习。最后，比较两种调优方法在各项自然语言处理任务上的学习曲线，并讨论其含义。

\subsection{结果与分析}
\subsubsection{文本分类}
首先，用传统的基于任务头的方法微调预训练语言模型，做文本分类。如第~3.1.2~节所述，这一过程用句子—标签对调优任务专用预测头，将语言模型（这里是 BERT）丰富的语言表示迁移到分类任务。数据集是 GLUE 中用于情感分析的 SST-2，由从电影评论中抽取的句子组成。每句话都由人工标注，表明它表达的是对电影的正面情感还是负面情感。把训练集分成正负两类，使用不同数量的样例训练模型，再在测试集上评估表现。

结果见表~\ref{tab:ch03-sst2} 的中列。零样本准确率接近 50\%，相当于随机猜测。在最初的几个样本规模上，准确率提高不多，到 64 个样例时有明显改善，最后在 256 个样例时达到 87\%。表中数值对实验具体选取哪些训练样例很敏感，多次运行所得的数值会略有不同，但总体趋势不变。

接着换用基于提示的训练进行比较。所得准确率见表~\ref{tab:ch03-sst2} 右列。每类仅用 32 个样例，提示模型的准确率就大幅上升。此后曲线很快趋于平缓，增加数据带来的收益较小。

\begin{table}[!htbp]
\centering
\caption{在 GLUE SST-2 数据集上，比较预训练再微调和基于提示的文本分类。表中给出不同训练样本数下的预测准确率。}
\label{tab:ch03-sst2}
\begin{tabular}{@{}rrr@{}}
\toprule
训练样本数 & PTFT 准确率 & 提示方法准确率\\
\midrule
0 & 0.5092 & 0.6800\\
16 & 0.5069 & 0.6743\\
32 & 0.6548 & 0.7867\\
64 & 0.8486 & 0.8475\\
128 & 0.8624 & 0.8521\\
256 & 0.8739 & 0.8658\\
\bottomrule
\end{tabular}
\end{table}

\sourcefigure{fig03_15.pdf}{表~\ref{tab:ch03-sst2} 的数据图。图中比较文本分类实验里预训练再微调与基于提示的学习的准确率随训练样例数的变化。}{fig:ch03-sst2}

图~\ref{fig:ch03-sst2} 从零样本出发，逐渐增加训练数据，绘出两种模型的准确率随训练样例数的变化。与基于任务头的分类器相比，提示模型只用极少的训练样例就取得了出色的结果。PTFT 模型最终可与基于提示的模型竞争，但要每类使用 32 个训练样例，才能达到与提示模型零样本表现相近的水平。\footnote{译注：原文称 32 个样例时已达到提示模型的零样本准确率。表~\ref{tab:ch03-sst2} 所列的两个数值分别为 0.6548 和 0.6800。}

\subsubsection{命名实体识别}
现在考察第二项分类任务，即命名实体识别。预计它在零样本条件下表现较差。我们要确定，在少样本条件下，提示方法是否仍优于微调。实验使用 CoNLL-2003 数据集，它收集了新闻标题和引语，逐词元标记人名、组织名、地名，另设一类容纳其他实体。

PTFT 的结果见表~\ref{tab:ch03-conll} 中列。与文本分类实验类似，BERT 在没有任务训练时表现较差。这并不令人意外，因为命名实体识别要识别标签，而 BERT 还不知道该数据集中各标签所指的含义。它预测的值大多大于 0，而多数标签等于 0，因而产生许多假阳性，$F_1$ 得分也较低。最初几个数据点显示，加入更多训练样例时，表现先变差，可能源于灾难性遗忘。随后，每增加一批数据，表现便有所改善，到 1024 个样例时达到 $F_1=85\%$。

\begin{table}[!htbp]
\centering
\caption{在 CoNLL-2003 数据集上，比较预训练再微调与基于提示的命名实体识别。表中给出不同训练样本数下的 $F_1$ 得分。}
\label{tab:ch03-conll}
\begin{tabular}{@{}rrr@{}}
\toprule
训练样本数 & PTFT 的 $F_1$ & 提示方法的 $F_1$\\
\midrule
0 & 0.0687 & 0.0712\\
8 & 0.0492 & 0.5788\\
16 & 0.0023 & 0.6482\\
32 & 0.0034 & 0.7274\\
64 & 0.3323 & 0.7867\\
128 & 0.5578 & 0.8365\\
256 & 0.7157 & 0.8672\\
512 & 0.7894 & 0.8304\\
1024 & 0.8526 & 0.8551\\
\bottomrule
\end{tabular}
\end{table}

接着采用基于提示的调优，逐次增加训练样例数，每次训练五轮，并查看学习曲线。表~\ref{tab:ch03-conll} 右列给出结果。仅用几句话，表现就大幅改善，并逐步提高，到 256 个样例时 $F_1$ 达到 87\%。使用更多数据调优也未能使表现进一步提高。

图~\ref{fig:ch03-conll} 汇总比较了 PTFT 和基于提示的调优。结果与文本分类类似。数据充足时，预训练再微调可与基于提示的学习竞争。数据缺乏时，基于提示的调优则取得了好得多的结果。

\sourcefigure{fig03_16.pdf}{表~\ref{tab:ch03-conll} 的数据图。图中比较命名实体识别实验里预训练再微调和基于提示的学习的 $F_1$ 得分随训练样例数的变化。}{fig:ch03-conll}

\subsection{结论}
本教程的两个实验共同表明，训练集规模有限时，基于提示的调优优于基于任务头的微调。取得规模足够大、质量良好的训练数据，长期以来都是许多机器学习问题的主要难点。因此，少样本学习取得的效果尤其值得注意。采用提示方法时，模型的质量取决于提示模板和答案格式的设计。选对这两项，只需几十个样例就能得到良好的自然语言处理结果。下一章将进一步讨论这一点。

\section{教程二——提示工程方法}
\label{sec:ch03-tutorial-two}
\subsection{概述}
本章的另一项主题是用模板工程提高经过提示调优的大语言模型的分析能力。清单~\ref{lst:ch03-capitals} 和表~\ref{tab:ch03-capitals} 说明，大语言模型的推理输出对模板结构及提示组成的细节很敏感。该演示使用了一个网页应用，适合验证思路，但因为不能微调，能力有限。本教程将通过少样本和多样本提示调优扩展这些实验，讨论不同提示模板的效果，说明优化提示模板的重要性。

\begin{keypoint}
\noindent\textbf{目标}
\begin{itemize}
\item 说明任务表现对提示模板设计十分敏感，即便细微改动，也能造成明显差异。
\item 考察哪些因素有助于提高提示模板的质量。
\item 用软提示进行自动调优，比较它与手工提示的效果。
\end{itemize}
\end{keypoint}

这里给出教程的简明版本。完整版见 \url{https://github.com/springer-llms-deep-dive/llms-deep-dive-tutorials}。

\subsection{实验设计}
本教程考察几种不同的模板工程方法，使用基准数据集训练模型并评估其表现。先用最简单的手工模板工程方法（见第~3.3.2~节）。这种方法由开发者构造最适合任务的模板。可以参考针对不同提示学习任务推荐模板的文献（例如第~3.2.4~节），也可以试验不同配置。

我们还会用训练数据探索自动模板设计。第~3.3.3 和 3.3.4~节已经介绍了几种自动方法。本教程采用一种基于梯度的模板优化方法，称为软提示。手工提示固定，软提示则使用可变模板，可以在不受离散词元映射限制的情况下将其调到更合适的表示。先用一个组合了数据集特征与“软”词元的模板初始化软提示，其中软词元可选择用特定的词或短语初始化。然后用反向传播调整相应嵌入，以达到训练目标。我们用 OpenPrompt 的代码实现软提示（\refcite{ch03-ding}{Ding et al., 2021）}。实验将说明，软提示可以优于手工构造的提示。

本教程用 SuperGLUE BoolQ 数据集。它提供由一段说明性文字、一条与段落内容有关的是非题和正确答案组成的三元组。BoolQ 的主题范围很广，包括历史、科学、地理、法律、体育和流行文化等，适合考察大语言模型的自然语言推断能力。

\subsection{结果与分析}
在手工提示实验中，我们用三种不同的提示模板和若干种训练样本数，对 \texttt{t5-base} 模型做基于提示的调优。训练样例在两个标签类别之间均分。我们把三种模板分别称为“最简单”模板、“最简单加标点”模板和更“合适”的模板。第一种只是把段落、问题和掩码连接起来。第二种用标点进一步指明文本结构。第三种加入文字提示各部分的含义，并把掩码词元放在形式合理的完形填空中。

\begin{lstlisting}[language=Python,nolol=true]
temp1 = "{passage} {question} {mask}"
temp2 = "{passage} . {question} ? {mask}"
temp3 = "hypothesis: {passage} premise: {question} The answer was {mask} ."
\end{lstlisting}

首先用最简单的模板做微调实验，结果见表~\ref{tab:ch03-boolq} 左侧的准确率列。总体来看，模型表现较差。

\begin{itemize}
\item 零样本推理将所有样本都判为负类，因此准确率等于验证集的负类样本比例（约为 50\%）。
\item 少样本结果有所改善，但只略好于随机猜测，相当于每道题抛一枚硬币。
\item 模型表现大约在 256 个样例时达到最高，准确率仍然较低。
\end{itemize}

\begin{table}[!htbp]
\centering
\caption{三种提示模板在 SuperGLUE BoolQ 上的零样本和少样本准确率。}
\label{tab:ch03-boolq}
\begin{tabular}{@{}rrrr@{}}
\toprule
训练样本数 & 最简单 & 最简单加标点 & 改进模板\\
\midrule
0 & 0.501 & 0.526 & 0.564\\
16 & 0.511 & 0.508 & 0.489\\
32 & 0.513 & 0.482 & 0.503\\
64 & 0.518 & 0.516 & 0.535\\
128 & 0.526 & 0.598 & 0.638\\
256 & 0.578 & 0.628 & 0.686\\
512 & 0.546 & 0.638 & 0.640\\
\bottomrule
\end{tabular}
\end{table}

接着试一项简单的改动，在段落末尾没有句号时加上句号，并在问题结尾加上问号。表~\ref{tab:ch03-boolq} 中间的准确率列显示了这个小改动的效果。只需在恰当的位置加入“?”和“.”，零样本准确率便由约 50\% 增至约 53\%。微调开始后，带标点模板的准确率上升得比最简单的模板快，说明提示有所改善。

最后试用更合适的手工模板。这个模板给出了有用的上下文，并明确向模型提问，预期效果更好。用改进后的模板重复相同实验，得到表~\ref{tab:ch03-boolq} 右列。图~\ref{fig:ch03-boolq} 给出三条学习曲线。改进模板的零样本表现明显优于另两种模板。不过，微调样例很少时，灾难性遗忘使准确率下降，继续微调之后才开始回升。使用 256 个样例后，模型能正确回答约 69\% 的提示，明显优于另两种模板。

\sourcefigure{fig03_17.pdf}{使用三种不同提示训练的模型在 SuperGLUE BoolQ 验证集上的准确率变化。模板形式会影响问答任务的零样本和少样本表现。}{fig:ch03-boolq}

接着用精心设计的 10 个候选模板重复上述实验，其中既有完形填空式，也有前缀式（详见完整版教程）。结果显示，模板有几个值得注意的特点。同一类模板的表现有一定共性。

\begin{itemize}
\item 前缀式提示在零样本条件下已有一定效果。用于调优的样例很少时，灾难性遗忘使表现下降。样本数进一步增加后，预测能力大幅提高。
\item 完形填空式提示的零样本表现稍差，只有少量训练样例时表现也会变差。随后，它的表现会改善，最终与前缀式提示相当。
\item 提供的上下文较少的提示，表现明显较差。零样本结果只略好于随机猜测，使用全部训练样例后，也只再提高几个百分点。不过，完全训练后，每条提示的表现都优于随机猜测，说明模型确实编码了其中一部分问题的答案。
\item 同一类中的不同模板，总体表现有明显差异，且随训练数据增加，差异往往也会增大。这说明模板结构中的细微差别也可能带来显著影响。
\end{itemize}

最后一个实验构造了两种软提示，一个采用简单模板，一个采用经过设计的模板。两种情况下，都用 128 个 SuperGLUE BoolQ 样例微调若干轮。这个实验冻结 \texttt{t5-base} 模型，只调优提示。图~\ref{fig:ch03-soft-boolq} 左图显示两种模型在验证集上的表现变化。

\begin{itemize}
\item 红色虚线代表缺少特征的提示。微调 60 轮后，表现仍无提高。这个模板的内容过少，缺少可以微调的特征，因此微调未能改善效果。
\item 黑色实线代表第二种提示。继续微调后，表现明显提高，60 轮内准确率约提高 4.5\%。与输入模板相比，软提示调优得到了更好的模板。
\end{itemize}

经过 60 轮训练，经过设计的软提示模型优于手工模板实验中所有模型的零样本表现，也优于这些模型用 128 个训练样本微调之前的表现。由于提示的规模远小于大语言模型本身，达到这一表现所需的算力较少。因此，当数据点有限、需要考虑训练算力开销时，提示调优可能优于大语言模型调优。

接着考察四种调优方法。

\begin{enumerate}
\item 仅调优软提示。
\item 仅调优大语言模型。
\item 先调优软提示，再调优大语言模型。
\item 同时调优软提示和大语言模型。
\end{enumerate}

图~\ref{fig:ch03-soft-boolq} 右图给出结果\footnote{译注：原图右侧横轴标为训练轮数，原文在下段按训练样例数讨论这些点。}，可以看到以下现象。

\begin{itemize}
\item 对于 BoolQ 任务，调优大语言模型明显优于单独调优软提示。只调提示的模型在继续调优后确实有明显改善，但在 512 个样例时，其表现远低于另外三个允许调整大语言模型参数的模型。随着训练时间延长，这一差距应会缩小。\refcite{ch03-lester}{Lester 等人（2021）}采用经过优化的训练方案，经过 30\,000 步得到约 0.9 的准确率。
\item 其他三个模型中，以调优过的软提示初始化的模型，零样本表现最好。它之前已经接触过训练数据，因此这一点并不出人意料。此后，允许调整大语言模型的几个模型，随训练样本增加而改善的速度相近。各个训练阶段，同时调优大语言模型和提示的模型略优于只调优大语言模型的模型。这很可能是因为两者共同调优时可调整的参数更多（2.48 亿对 2.22 亿），它也是本教程中表现最好的模型。不过，与使用经过设计的提示并调优大语言模型的做法相比，提升幅度很小。
\end{itemize}

\sourcefigure{fig03_18.pdf}{左图是从简单提示和经过设计的提示出发，进行软提示调优的结果。右图是图例所列四种学习方式的学习曲线。同时调优提示和大语言模型的模型，在训练过程的各个阶段表现最好。}{fig:ch03-soft-boolq}

\subsection{结论}
实验表明，提示工程对改善大语言模型表现很重要。要取得最佳表现，当然还有许多可调因素，这里没有重点考察，例如训练集规模、训练轮数、学习率和大语言模型的选择等。尽管如此，从表现最差的模型到表现最好的模型，仅通过改变模板设计，预测准确率就提高了 25\% 以上。因此，使用提示的模型需要认真考虑模板设计。
